\section{习题}

\subsection{基础过关题}

\begin{enumerate}

\item \textbf{（单项选择题）} 下列关于操作系统的叙述中，正确的是
\begin{center}
\begin{tabularx}{\linewidth}{@{}XX@{}}
(A) 操作系统是应用软件，与用户程序地位相同 & (B) 操作系统管理处理器、存储器、I/O 设备和文件，并为用户和应用程序提供接口 \\[4pt]
(C) 并发是指同一时刻多个程序同时执行 & (D) 虚拟只能把多个物理实体合并为一个逻辑实体
\end{tabularx}
\end{center}

\ifshowsolutions\par\noindent\textbf{解析：} 操作系统是最基本的\emph{系统软件}，它运行在内核态，用户程序运行在用户态并受其管理，二者地位并不相同，故 (A) 错。操作系统的功能可以概括为两类：一是资源管理，即对处理机、存储器、I/O 设备和文件进行分配、回收与调度；二是提供接口，即通过命令接口、图形界面和系统调用为使用者与应用程序服务，(B) 正确。并发是「宏观上同时、微观上交替」，单核 CPU 用时间片轮转就能实现；同一时刻真正同时执行多个指令流称为并行，必须有多个处理器核，(C) 把并发说成了并行。虚拟的常见形式恰恰是「一虚多」：一个物理 CPU 虚拟成多个逻辑 CPU、虚拟存储器让每个进程都以为自己独占内存、SPOOLing 把一台打印机虚拟成多台；(D) 说「只能」多虚一，错。补充一句：共享分为互斥共享（打印机等临界资源）和同时访问（磁盘、只读代码段），并非所有资源都能被同时共享。\boxed{\text{答案：(B)}}\else \answerarea \fi

\vspace{8pt}

\item \textbf{（单项选择题）} 下列事件中，不可能引起 CPU 从用户态转入内核态的是
\begin{center}
\begin{tabular}{ll}
(A) 执行除法指令时除数为 $0$ & (B) 用户程序调用 \texttt{read()} 读文件 \\[4pt]
(C) 时钟部件发出时钟中断 & (D) 执行一条把两个寄存器内容相加的指令
\end{tabular}
\end{center}

\ifshowsolutions\par\noindent\textbf{解析：} 用户态到内核态的转换只有三条途径：中断、异常和系统调用（陷入）。(A) 除零是 CPU 执行指令时发现的同步异常，CPU 会转到异常处理程序，进入内核态。(B) \texttt{read()} 是库函数，其内部通过陷入指令发起系统调用，CPU 由用户态进入内核态执行服务例程。(C) 时钟中断是外部中断，CPU 在执行完当前指令后响应，转入内核态的中断服务程序——它也是分时与抢占式调度的基础。(D) 寄存器相加只涉及用户可见的寄存器和运算器，不访存、不请求服务、也不触发异常，CPU 始终停留在用户态。\boxed{\text{答案：(D)}}\else \answerarea \fi

\vspace{8pt}

\item \textbf{（单项选择题）} 下列指令中，只能在内核态执行的是
\begin{center}
\begin{tabular}{ll}
(A) 关中断（清除 PSW 中的中断允许位） & (B) 把两个通用寄存器的内容相加 \\[4pt]
(C) 读取本进程的程序计数器 & (D) 转移到本进程代码段内的另一地址
\end{tabular}
\end{center}

\ifshowsolutions\par\noindent\textbf{解析：} 关中断会让 CPU 不再响应时钟等外部中断，操作系统将失去对时间片和并发的控制，因此它是\emph{特权指令}，只能在内核态执行，(A) 正确。(B) 是普通的算术运算，(C) 只读本进程可见的 PC，(D) 是段内跳转，三者都在本进程的地址空间与寄存器范围内操作，不会影响其他进程和系统状态，属于用户态可执行的普通指令。同类的特权指令还有：I/O 指令、修改页表基址寄存器、设置时钟/定时器、修改 PSW、停机指令。判断标准只有一条：该指令是否会改变系统全局状态或访问受保护的硬件资源。\boxed{\text{答案：(A)}}\else \answerarea \fi

\vspace{8pt}

\item \textbf{（填空题）} 用户程序执行陷入指令请求系统调用，按异常分类属于 \underline{\hspace{2.2cm}} \ \textcircled{1}，处理结束后返回到陷入指令的 \underline{\hspace{2.2cm}} \ \textcircled{2}。缺页由访存指令触发，按异常分类属于 \underline{\hspace{2.2cm}} \ \textcircled{3}，缺页处理完成后 CPU \underline{\hspace{2.2cm}} \ \textcircled{4} 引发缺页的那条指令。

\ifshowsolutions\par\noindent\textbf{解析：} 异常按「处理完后从哪里继续」分三类。陷阱（Trap）是程序\emph{故意}执行一条指令造成的，如系统调用、断点，它不算错误，因此处理完返回下一条指令（陷入指令本身已经执行完）。故障（Fault）是当前指令执行\emph{失败但可恢复}，如缺页、除零（能否恢复依体系结构而定），处理完必须\emph{重新执行}出错的那条指令，否则这次访存就丢了。终止（Abort）是不可恢复的严重错误，如硬件校验错，通常直接终止进程，没有返回点。还要区分：中断来自 CPU 外部，与当前指令异步；异常由当前指令触发，与指令同步。\boxed{\text{答案：\textcircled{1}陷阱（Trap）；\textcircled{2}下一条指令；\textcircled{3}故障（Fault）；\textcircled{4}重新执行}}\else \answerarea \fi

\vspace{8pt}

\item \textbf{（单项选择题）} 关于微内核结构，下列说法正确的是
\begin{center}
\begin{tabularx}{\linewidth}{@{}XX@{}}
(A) 它把文件系统、设备驱动等主要服务放在内核态，以减少态切换 & (B) 它只在内核保留调度、进程间通信、中断处理等基本机制，可靠性较高，但态切换开销较大 \\[4pt]
(C) 它使系统不再发生用户态与内核态的切换 & (D) 它允许两个用户态服务进程直接调用对方的函数，无需经过内核
\end{tabularx}
\end{center}

\ifshowsolutions\par\noindent\textbf{解析：} (A) 描述的其实是大内核：功能模块都在内核态，模块间是普通函数调用，开销小，但内核庞大、一处出错可能拖垮整个系统。(B) 正确：微内核只保留进程调度、进程间通信、中断处理、基本地址空间管理等最小机制，文件系统、设备驱动、网络协议栈等以用户态服务进程的形式运行，单个服务崩溃不会波及内核，也便于按需增删服务；代价是每次服务请求都要「陷入内核—由内核转发消息—再返回」，态切换与消息传递次数明显增加。(C) 错：微内核服务仍在用户态，每次请求仍要陷入内核，切换反而更频繁。(D) 错：两个用户态服务位于各自独立的地址空间，不能直接互相调用，必须通过内核的消息传递（IPC）通信。\boxed{\text{答案：(B)}}\else \answerarea \fi

\vspace{8pt}

\item \textbf{（填空题）} 某处理器主频为 $2$ GHz。程序 $P$ 执行 $8\times 10^{6}$ 条用户态指令，平均 CPI 为 $1$；执行过程中共发起 $1000$ 次系统调用，每次从陷入到返回需要 $2000$ 个时钟周期（已包含态切换、参数检查与服务例程本身，忽略其他开销）。$P$ 的执行时间为 \underline{\hspace{2.2cm}}，系统调用开销所占比例为 \underline{\hspace{2.2cm}}。

\ifshowsolutions\par\noindent\textbf{解析：} 两类开销分别折算成时钟周期再相加。用户态指令：$8\times10^{6}\times1=8\times10^{6}$ 个周期；系统调用：$1000\times2000=2\times10^{6}$ 个周期。总计 $1.0\times10^{7}$ 个周期。主频 $2$ GHz 即每秒 $2\times10^{9}$ 个周期，所以
\[
T=\frac{1.0\times10^{7}}{2\times10^{9}}\ \text{s}=5\times10^{-3}\ \text{s}=5\ \text{ms}.
\]
其中系统调用部分为 $2\times10^{6}/2\times10^{9}=1$ ms，占 $1/5=20\%$；平均每次系统调用耗时 $2000/(2\times10^{9})\ \text{s}=1\ \mu\text{s}$。注意两点：用户指令的 CPI 只对用户态指令有效，系统调用必须按「次数 $\times$ 每次周期数」单独累加，不能套用同一个 CPI；百分比是「系统调用周期 $\div$ 总周期」，不是 $\div$ 用户指令周期。\boxed{\text{答案：$5$ ms；$20\%$}}\else \answerarea[4] \fi

\end{enumerate}

\subsection{真题风格与综合训练}

\begin{enumerate}[resume]

\item \textbf{（真题风格·综合题）} 某单 CPU 系统采用多道程序设计，内存可同时容纳 $3$ 道作业，忽略调度与进程切换开销。三道作业同时到达，各自依次执行「CPU 段—I/O 段—CPU 段」，时间如下（单位：ms）。

\begin{center}
\begin{tabular}{|c|c|c|c|}
\hline
作业 & 第 $1$ 段 CPU & I/O & 第 $2$ 段 CPU \\
\hline
J1 & $20$ & $40$ & $10$ \\
J2 & $10$ & $30$ & $20$ \\
J3 & $20$ & $20$ & $20$ \\
\hline
\end{tabular}
\end{center}

假设：各作业的 I/O 由独立设备完成，可彼此并行，也可与 CPU 重叠执行；CPU 采用非抢占调度，空闲时若有多道作业就绪，按 J1、J2、J3 的固定次序选择。

（1）若按 J1$\to$J2$\to$J3 单道顺序执行，总时间是多少？
（2）多道执行时，画出 CPU 时间线，求总完成时间与 CPU 利用率（保留三位有效数字）。
（3）求多道执行时各作业的周转时间。

\ifshowsolutions\par\noindent\textbf{解析：}
（1）单道执行时同一时刻只有一个作业在系统中，CPU 在它做 I/O 期间只能空闲。三个作业的 CPU 时间与 I/O 时间完全串行：
\[
(20+40+10)+(10+30+20)+(20+20+20)=70+60+60=190\ \text{ms}.
\]

（2）按时刻逐个推进。$0$ 时三道作业都到达，选 J1：$0\sim20$ 运行 J1 第 $1$ 段；$20$ 时 J1 转入 I/O（$20\sim60$），CPU 交给 J2，$20\sim30$ 运行 J2 第 $1$ 段；$30$ 时 J2 转入 I/O（$30\sim60$），CPU 交给 J3，$30\sim50$ 运行 J3 第 $1$ 段；$50$ 时 J3 转入 I/O（$50\sim70$），而 J1、J2 的 I/O 要到 $60$ 才完成，故 CPU 在 $50\sim60$ 空闲；$60$ 时 J1、J2 同时就绪，按题设次序选 J1，$60\sim70$ 运行 J1 第 $2$ 段；$70$ 时 J1 结束，J3 的 I/O 也恰好完成，就绪的有 J2、J3，选 J2，$70\sim90$；最后 $90\sim110$ 运行 J3 第 $2$ 段。

\begin{center}
{\small
\setlength{\tabcolsep}{4pt}%
\begin{tabular}{|c|c|c|c|c|c|c|c|}
\hline
时段/ms & 0--20 & 20--30 & 30--50 & 50--60 & 60--70 & 70--90 & 90--110 \\
\hline
CPU & J1 & J2 & J3 & 空闲 & J1 & J2 & J3 \\
\hline
\end{tabular}}
\end{center}

总完成时间为 $110$ ms。CPU 忙碌时间为 $20+10+20+10+20+20=100$ ms，故利用率为
\[
\frac{100}{110}\approx 0.909,
\]
即 $90.9\%$。作为对照，单道执行的利用率为 $100/190\approx52.6\%$：多道的作用正是用另一道作业填补 I/O 造成的 CPU 空档，本题中空闲时间由 $90$ ms 压缩到 $10$ ms。

（3）三道作业都在 $0$ 时刻到达，周转时间就等于完成时刻：J1 为 $70$ ms，J2 为 $90$ ms，J3 为 $110$ ms，平均 $(70+90+110)/3=90$ ms。

\boxed{\text{答案：（1）$190$ ms；（2）$110$ ms、$90.9\%$；（3）$70$、$90$、$110$ ms（均值 $90$）}}\else \answerarea[10] \fi

\vspace{8pt}

\item \textbf{（真题风格·单项选择题）} 单 CPU 系统中，用户进程 A 正在用户态执行，此时磁盘控制器发出一个 I/O 完成中断，该 I/O 是此前进程 B 请求的。下列叙述中一定成立的是
\begin{center}
\begin{tabularx}{\linewidth}{@{}XX@{}}
(A) 进程 A 立即被置为阻塞态 & (B) CPU 执行完当前指令后由用户态转入内核态，执行中断服务程序 \\[4pt]
(C) 操作系统立即切换到进程 B 运行 & (D) 中断处理结束后，CPU 一定返回 A 的断点继续执行
\end{tabularx}
\end{center}

\ifshowsolutions\par\noindent\textbf{解析：} 外部中断与当前指令异步，CPU 在一条指令执行结束时检查中断请求，若有未被屏蔽的中断，就保存断点（PC）和 PSW，按中断向量转入内核态的中断服务程序，(B) 一定成立。(A) 错：这个 I/O 属于 B，与 A 无关，A 不会因此被阻塞；A 最多在返回用户态时被调度程序换下，那是调度决策而不是「立即阻塞」。(C) 错：中断服务程序只做必要的工作（标记 I/O 完成、把 B 唤醒成就绪态），是否切换、切换到谁由调度程序决定，且切换一般发生在中断处理结束、即将返回用户态的时刻。(D) 错：返回时若调度程序选中了就绪的其他进程，CPU 会转到那个进程，并不保证回到 A。\boxed{\text{答案：(B)}}\else \answerarea \fi

\vspace{8pt}

\item \textbf{（真题风格·综合题）} 用户进程 A 调用 \texttt{read(fd, buf, 100)} 从磁盘文件读取 $100$ 字节，所读数据不在页缓存中，需要启动磁盘 I/O。题述事件如下：

\textcircled{1} CPU 执行陷入指令，把断点（PC）与 PSW 等现场保存到内核栈，切换到内核态；\\
\textcircled{2} 内核按系统调用号查系统调用表，找到 \texttt{read} 的服务例程；\\
\textcircled{3} 内核检查参数后启动磁盘读，把 A 置为阻塞态，调度进程 B 运行；\\
\textcircled{4} 磁盘控制器在数据就绪后发出 I/O 完成中断；\\
\textcircled{5} 中断服务程序把数据从磁盘缓冲区送入页缓存，并把 A 唤醒；\\
\textcircled{6} 内核把数据从页缓存拷入 A 的用户缓冲区，恢复 A 的现场，返回用户态，\texttt{read} 返回实际读到的字节数。

假设系统在中断处理结束后立即调度 A 运行，且进程切换本身在内核态内部完成、不额外产生用户态与内核态之间的切换。

（1）按发生顺序排列上述事件。
（2）从 A 发起 \texttt{read} 到 A 收到返回值，A 自身经历了几次「用户态$\to$内核态」、几次「内核态$\to$用户态」？整个系统（含中断）共几次？
（3）为什么 A 不能直接调用内核中的 \texttt{read} 服务例程，而必须通过系统调用？

\ifshowsolutions\par\noindent\textbf{解析：}
（1）顺序为 \textcircled{1}$\to$\textcircled{2}$\to$\textcircled{3}$\to$\textcircled{4}$\to$\textcircled{5}$\to$\textcircled{6}。先是由用户态陷入内核（\textcircled{1}），再在内核中定位服务例程（\textcircled{2}）；参数合法但数据不在内存，于是发起磁盘读并让出 CPU（\textcircled{3}）；磁盘完成后由中断驱动后续工作（\textcircled{4}）；中断服务程序完成数据传送并唤醒 A（\textcircled{5}）；A 再次被调度时，内核替它完成系统调用的返回（\textcircled{6}）。

（2）计数口径要分清「A 自身」和「整个系统」。
对 A 自身：只在 \textcircled{1} 发生一次用户态$\to$内核态（陷入），只在 \textcircled{6} 发生一次内核态$\to$用户态（系统调用返回），共 $2$ 次。
对整个系统：用户态$\to$内核态共 $2$ 次——\textcircled{1} 的陷入，以及 \textcircled{4} 的磁盘中断（当时运行的 B 从用户态进入内核）；内核态$\to$用户态也共 $2$ 次——\textcircled{3} 调度 B 之后内核必须返回 B 的用户态，\textcircled{6} 返回 A 的用户态。合计 $4$ 次。注意 \textcircled{5} 在中断处理内部完成，本身不产生态切换；\textcircled{3} 的进程切换在内核态内部完成，也不额外产生「用户态/内核态」的切换——切换的是进程，不是特权级。

（3）内核的代码与数据位于内核地址空间，页表/段表中这些页被标记为仅内核可访问（如 x86 页表项的 U/S 位），CPU 在用户态处于低特权级，访问它们会触发保护性异常，而不会执行内核代码；修改页表基址寄存器、开关中断等特权指令在用户态同样非法。因此用户程序无法用普通 \texttt{call} 进入内核。系统调用提供唯一的受控入口：用体系结构规定的陷入指令把控制权交给内核中预先注册的入口，用系统调用号选择具体服务，内核在服务前还会检查参数（如用户缓冲区指针是否落在该进程的用户空间内），既完成服务又不破坏隔离。\boxed{\text{答案：（1）\textcircled{1}$\to$\textcircled{2}$\to$\textcircled{3}$\to$\textcircled{4}$\to$\textcircled{5}$\to$\textcircled{6}；（2）A：$2$ 次，系统：$4$ 次；（3）内核页受保护}}\else \answerarea[10] \fi

\vspace{8pt}

\item \textbf{（真题风格·单项选择题）} 下列关于操作系统体系结构的叙述中，正确的是
\begin{center}
\begin{tabularx}{\linewidth}{@{}XX@{}}
(A) 层次结构中每一层只使用其相邻下层提供的功能，因此可自底向上逐层调试 & (B) 模块化结构中各模块通过接口相互调用，因此模块之间不存在依赖关系 \\[4pt]
(C) 简单结构（如 MS-DOS）划分了清晰的模块边界，用户程序不能直接访问硬件 & (D) 外核把进程调度与内存分配的具体策略都实现在内核中，应用只能被动接受
\end{tabularx}
\end{center}

\ifshowsolutions\par\noindent\textbf{解析：} 层次结构把操作系统分成若干层，每层只能调用相邻下层提供的接口，层间是单向依赖；正因为如此，实现和调试都可以自底向上进行，调试第 $k$ 层时只需信任已调试好的下面各层，(A) 正确。(B) 错：模块化结构按功能划分模块并规定接口，模块之间通过接口调用，调用关系可以交叉，因而存在相互依赖，这也使得模块化结构难以做到严格的顺序调试。(C) 错：简单结构（MS-DOS）没有清晰的模块划分，应用甚至能越过操作系统直接访问硬件，几乎没有保护。(D) 错：外核（Exokernel）只在内核中负责资源的保护与多路复用，把具体的资源管理策略（调度策略、页面置换策略等）下放给应用层的库操作系统，说反了。\boxed{\text{答案：(A)}}\else \answerarea \fi

\end{enumerate}

\vspace{8pt}
